\documentclass{article}
\begin{document}
\begin{abstract}
Noisy supervision makes a fine-tuned model memorize wrong labels.

We filter samples by gradient alignment and gain 2.1 points.
\end{abstract}

\section{Introduction}
Retrieval-augmented generation (RAG) and bidirectional encoders (BERT) are standard. RAG is used here once.
We evaluate on HumanEval with HuggingFace tooling and call the method GradSieve.
The remainder of this paper is organized as follows.

\section{Method}
Moreover, the filter keeps aligned samples. Furthermore, it drops the rest. Additionally, it is cheap.
The filter keeps a sample; the rest are dropped.
It is worth noting that the threshold matters.
The result may potentially hold to some extent, and it might arguably be relatively stable.
The rule has two parts: a threshold and a reference.
The reference gradient --- computed on clean data --- is reused.
The \textbf{threshold} is fixed.
A myriad of settings were tried and the gain was +2.1\% on average.
The gradient-alignment-based filter is simple.

The filter compares each sample with the reference and decides which samples stay so that the model
can focus on trustworthy samples during the whole process and the procedure becomes more stable overall.

In summary, the rule depends on one threshold and is applied inside every batch before the update step.

\begin{proof}
We first note that the retained set has a mean gradient whose cosine with the reference is bounded below
by the threshold because every retained sample satisfies the inequality and the mean of vectors that each
satisfy a linear inequality satisfies the same inequality after normalization by the triangle inequality
and then we apply the descent lemma with the stated step size to conclude that the reference loss goes down
by the claimed amount which depends on the product of the two norms and the constant kappa defined above
so the claim follows for all steps of the procedure as long as the assumptions of the theorem hold and the
relative error of the reference gradient stays below the threshold which is the design rule of the method
and therefore nothing more needs to be shown here for the statement of the theorem to be established.
\end{proof}
Several works \cite{a,b,c,d,e} explored this.

\section{Conclusion}
In summary, the rule works.
\end{document}
